% Предложения по склейкам (нарративным переходам) между разделами.
% Не входит в сборку диссертации. Компилируется автономно (xelatex).
\documentclass[12pt,a4paper]{article}
\usepackage{fontspec}
\setmainfont{Times New Roman}
\setmonofont{Consolas}
\usepackage{polyglossia}
\setdefaultlanguage{russian}
\setotherlanguage{english}
\usepackage[margin=2.3cm]{geometry}
\usepackage{amsmath}
\usepackage[dvipsnames]{xcolor}
\usepackage{mdframed}
\usepackage{enumitem}
\usepackage{hyperref}
\hypersetup{colorlinks=true, linkcolor=blue, urlcolor=blue}

% блок для существующего наброска \link
\newenvironment{sketch}{%
  \begin{mdframed}[backgroundcolor=black!4, linecolor=black!25, linewidth=0.4pt,
    innertopmargin=5pt, innerbottommargin=5pt]\small\itshape
  \textbf{\upshape Набросок \texttt{\textbackslash link}:}\ }{%
  \end{mdframed}}
% блок для предлагаемой склейки
\newenvironment{glue}{%
  \begin{mdframed}[backgroundcolor=green!5, linecolor=green!40!black, linewidth=0.6pt,
    innertopmargin=5pt, innerbottommargin=5pt]
  \textbf{Предлагаемая склейка:}\par\smallskip}{%
  \end{mdframed}}

\newcommand{\loc}[1]{\par\noindent\textcolor{blue!55!black}{\textbf{Где:} #1}\par}
\newcommand{\note}[1]{\par\noindent\textcolor{red!60!black}{\textbf{Заметка:} #1}\par}

\title{Предложения по склейкам (нарративным переходам)\\между разделами диссертации}
\author{}
\date{\today}

\begin{document}
\maketitle

\noindent\textit{Документ собирает существующие рабочие пометки \texttt{\textbackslash link\{\}}
и предлагает для каждой готовую формулировку перехода. Зелёные блоки —
текст-кандидат на замену \texttt{\textbackslash link} в основном тексте; серые —
исходные наброски; красным — структурные замечания. Документ не входит в сборку
диссертации.}

\vspace{0.5em}\hrule\vspace{0.8em}

\section*{Обзор}
В собираемом тексте семь рабочих связок:
\begin{enumerate}[nosep]
  \item Гл.~1: биологические сети $\to$ искусственные (перед \S1.5).
  \item Гл.~1: анализ представлений $\to$ синтез стимулов (перед \S1.6).
  \item Гл.~2: геометрия гиппокампа $\to$ РНС (перед разделом о РНС).
  \item Гл.~2: завершение популяционного уровня $\to$ индивидуальный (хвост гл.~2).
  \item Гл.~3: INTENSE $\to$ вовлечённость в популяционный код (в разделе DRIADA).
  \item Раздел о вовлечённости: внутренний мост индивидуальный$\leftrightarrow$популяционный.
  \item Гл.~3: корреляционные методы $\to$ синтез стимулов (перед разделом MANGO).
\end{enumerate}
Связки 1, 2, 7 уже близки к финальному виду; 3 и 6 требуют лёгкой полировки;
4 и 5 содержат структурные проблемы (отмечены).

\section{Гл.~1: биологические $\to$ искусственные сети}
\loc{Перед \S\,1.5 «Интерпретация внутренних представлений в ИНС» (\texttt{draft\_ann\_repr}).}
\begin{sketch}
Предыдущие секции были посвящены методам анализа активности биологических
нейронных сетей. Однако задача интерпретации внутренних представлений актуальна
и для искусственных нейронных сетей. Более того, глубокие нейронные сети
демонстрируют замечательные параллели с биологическими системами\ldots
\end{sketch}
\begin{glue}
Рассмотренные выше методы~--- снижение размерности, оценка размерности, сетевой
анализ, анализ селективности~--- разрабатывались для активности биологических
нейронных сетей. Однако задача интерпретации внутренних представлений в равной
мере актуальна и для искусственных сетей, а глубокие сети обнаруживают глубокие
параллели с мозгом в организации кода~--- от селективности отдельных элементов
до низкоразмерной структуры популяционной активности. Эти параллели делают
перенос аналитического инструментария между нейронаукой и машинным обучением не
только возможным, но и продуктивным; ему посвящён настоящий раздел.
\end{glue}

\section{Гл.~1: анализ представлений $\to$ синтез стимулов}
\loc{Перед \S\,1.6 «Методы синтеза оптимальных стимулов» (\texttt{draft\_stimuli}).}
\begin{sketch}
В предыдущей секции обсуждались методы анализа уже сформированных представлений
\ldots Однако задачу интерпретации можно поставить принципиально иначе: не
анализировать отклик на заданные стимулы, а \textit{синтезировать} стимул\ldots
\end{sketch}
\begin{glue}
Все перечисленные подходы~--- и для биологических, и для искусственных сетей~---
анализируют уже сформированные представления: они описывают отклик нейрона на
заданные стимулы. Задачу интерпретации, однако, можно обратить: не анализировать
реакцию на предъявленное, а \textit{синтезировать} стимул, максимально
активирующий данный элемент. Такой подход~--- \textbf{максимизация активации}~---
переводит исследование из режима пассивного наблюдения в режим активного
зондирования репрезентаций и применим как к искусственным, так и к биологическим
нейронам.
\end{glue}

\section{Гл.~2: геометрия гиппокампа $\to$ рекуррентная сеть}
\loc{Перед разделом о РНС (\texttt{draft\_rnn}, \S\,2.3).}
\begin{sketch}
Коллективные переменные возникают и в искусственных нейронных сетях, особенно
в тех, где внутренняя динамика достаточно богата для формирования аттракторов.
Мы проверяли, можно ли из коллективной активности РНС сделать вывод о механизмах
принятия решения\ldots
\end{sketch}
\begin{glue}
Восстановление геометрии среды из активности гиппокампа показало, что коллективные
переменные несут структуру, изоморфную пространству задачи. Такие низкоразмерные
коллективные переменные возникают и в искусственных сетях с достаточно богатой
внутренней динамикой~--- прежде всего рекуррентных, способных формировать
аттракторы. В следующем разделе те же методы снижения размерности применяются
к скрытой активности рекуррентной сети, чтобы выяснить, можно ли по её
коллективной динамике судить о механизме принятия решения и о том, как алгоритм
обучения формирует вычислительную стратегию.
\end{glue}

\section{Гл.~2: завершение популяционного уровня $\to$ индивидуальный}
\loc{Хвост гл.~2 (\texttt{part2.tex:202}), после разделов о размерности и
мультиплексном графе.}
\begin{sketch}
Нужно специальное решение для ситуаций, когда одна или несколько популяционных
переменных «забивают» сигнал от всех остальных, как это происходит с кодированием
места в гиппокампе.
\end{sketch}
\note{Сейчас эта связка «висячая»: раздел-адресат (о гетерогенности/вовлечённости,
\texttt{popcode}) не подключён и физически находится в гл.~3. Предлагается
переформулировать её как \emph{переход между главами} 2$\to$3.}
\begin{glue}
Описанные методы~--- снижение размерности, оценка размерности, мультиплексный
граф рекуррентности~--- характеризуют популяцию как единое целое. Они, однако,
не разделяют вклад отдельных нейронов в коллективные переменные, а в гиппокампе
ситуация осложняется тем, что пространственное кодирование доминирует и «забивает»
прочие коллективные моды. Чтобы связать популяционный и индивидуальный уровни~---
выяснить, какие именно нейроны формируют ту или иную коллективную переменную,~---
необходим анализ индивидуальной селективности, которому посвящена следующая глава.
\end{glue}

\section{Гл.~3: INTENSE $\to$ вовлечённость в популяционный код}
\loc{В разделе DRIADA (\texttt{part3.tex:10}); адресат~--- раздел о вовлечённости
(\texttt{draft\_population\_involvement}).}
\begin{sketch}
INTENSE можно использовать не только для индивидуальной нейронной селективности,
но и для анализа вовлечённости тех или иных нейронов в популяционные переменные,
тем самым исследовать их нейронную основу.
\end{sketch}
\note{Раздел-адресат пока не подключён (\texttt{\textbackslash input} отсутствует);
рядом стоит устаревший \texttt{\textbackslash source} с числом 25\,\%, которое
новый расчёт на 64 сессиях не подтверждает (\,$\approx$49\,\%, обогащение
$\approx$1.1$\times$). При интеграции обновить число.}
\begin{glue}
Метод INTENSE измеряет селективность нейрона к внешним переменным. Тот же
информационно-теоретический критерий применим и к коллективным переменным,
извлечённым из популяционной активности: рассматривая каждую коллективную
переменную как ещё одну динамическую величину, можно определить, какие нейроны
её формируют. Это замыкает индивидуальный и популяционный уровни анализа в единой
схеме~--- анализу вовлечённости нейронов в популяционный код посвящён следующий
раздел.
\end{glue}

\section{Раздел о вовлечённости: внутренний мост}
\loc{Открывающая связка раздела \texttt{draft\_population\_involvement} (уже написана).}
\note{Этот \texttt{\textbackslash link} уже сформулирован как полноценный переход
и пригоден без правок; приведён для полноты.}
\begin{sketch}
Предыдущие главы рассматривали два уровня описания нейронной активности по
отдельности~--- популяционный и индивидуальный. Оба уровня опираются на одни и
те же экспериментальные записи, поэтому естественно поставить вопрос об их связи
\ldots Метод INTENSE применим без изменений и к этой задаче.
\end{sketch}

\section{Гл.~3: корреляционные методы $\to$ синтез стимулов (MANGO)}
\loc{Перед разделом MANGO (\texttt{draft\_mango\_intro:5}).}
\begin{sketch}
Методы, описанные в предыдущих секциях, анализируют корреляции между
зарегистрированной активностью нейронов и наблюдаемым поведением\ldots Однако
задачу интерпретации функции нейрона можно поставить принципиально иначе:
\textit{создать} оптимальный стимул, вызывающий максимальный отклик\ldots
\end{sketch}
\note{В \texttt{part3.tex:78} уже есть вводный абзац перед разделом MANGO с тем же
смыслом. Чтобы не дублировать, стоит оставить \emph{один} переход~--- предлагаемый
ниже~--- и убрать второй.}
\begin{glue}
INTENSE, снижение размерности и сетевой анализ выявляют статистические зависимости
между активностью нейрона и уже зарегистрированными переменными. Функцию нейрона
можно, однако, исследовать и активно: не искать корреляции с наблюдаемым, а
синтезировать стимул, вызывающий максимальный отклик. Такой подход~--- максимизация
активации~--- дополняет корреляционный анализ со стороны генерации и применим к
нейронам сетей любой природы; его реализации для импульсных сетей посвящён
настоящий раздел.
\end{glue}

\section*{Дополнительно: межглавные переходы}
Рабочих набросков для переходов между главами нет, но они~--- наибольший вклад
в связность (см. рецензию). Предложения, выводимые из общей идеи (квадрат
«биология/ИНС $\times$ популяция/индивид»):

\loc{Конец гл.~1 (обзор) $\to$ начало гл.~2.}
\begin{glue}
Обзор показал, что методы анализа коллективной и индивидуальной активности
развиты по отдельности и редко связаны между собой, а их перенос на искусственные
сети остаётся фрагментарным. Дальнейшие главы восполняют это: глава~2 посвящена
популяционному уровню~--- восстановлению и количественной характеризации скрытых
переменных коллективной динамики, глава~3~--- индивидуальной селективности и
инструментам, связывающим оба уровня. Обе главы последовательно проходят путь
от биологических сетей к искусственным.
\end{glue}

\loc{Конец гл.~2 $\to$ начало гл.~3.}
\begin{glue}
См. склейку №\,4: завершив популяционное описание, переходим к индивидуальному
уровню и инструментам, замыкающим оба масштаба.
\end{glue}

\section*{Сводка действий}
\begin{enumerate}[nosep]
  \item Склейки 1, 2, 7~--- внести как финальный текст (минимальная полировка).
  \item Склейка 3~--- внести с привязкой к результату о геометрии гиппокампа.
  \item Склейка 4~--- переформулировать висячую связку \texttt{part2:202} как
        переход 2$\to$3.
  \item Склейка 5~--- внести \emph{после} интеграции раздела о вовлечённости и
        обновления числа 25\,\%$\to$новый вывод.
  \item Склейка 7~--- устранить дублирование с \texttt{part3:78}.
  \item Добавить два межглавных перехода (1$\to$2, 2$\to$3).
\end{enumerate}

\end{document}
